\quellseite{20}
\part{Part I. Introductory}

\quellseite{22}
\chapter[Chapter I. The Catechist’s Office and Work]{The Lutheran Catechist.\\Chapter I. The Catechist’s Office and Work}

In the whole animal kingdom there is probably not a single species that comes into the world so helpless as man. He enters upon life a weak, dependent, ignorant and depraved being.\randnotiz{Man’s need of instruction.} Left to himself he must inevitably and miserably perish. And perish he does wherever he is thus abandoned. The first care he requires is that his physical life may remain in him and develop. For years he requires attention, care and nurture, that he may gradually grow into self-help and independence.

This is true of his mind as well as of his body. Left entirely to self, without contact with other and superior minds, his mind would remain practically a blank, and would rise little, if any, above the instincts of the brute. If he came in contact with other minds filled with savagery and superstition he would imbibe a certain kind and amount of knowledge. He might gather a store of empirical shrewdness and wisdom. He would acquire little, if any, of
\quellseite{23}
the refining and ennobling truth that he so much needs. He would pick up much that is false. He would accept many superstitions as truth. He would accept dangerous beliefs and fall into unworthy and hurtful practices. He needs a teacher and a guide to direct and develop his mind aright. He needs instruction in that which is true and good.

The mind of man has a spiritual side. The spiritual nature of man was made for God; it is out of its true element, unsatisfied and undeveloped until it rests in God. Draped in darkness and shrouded in sin he is “but a child crying in the night; only a child, crying for a light.” Reaching up and crying out, he seeks for satisfying truth. Without a teacher he finds at best but fragments of truth mixed with much degrading and destructive error. He needs a spiritual teacher, a spiritual guide to give to him the truth he needs.

\randnotiz{The first teacher.} In Christian lands the first teacher ought to be the mother. She is the first, the God-given teacher of the child. The father is to stand by and assist the mother all he can. We take for granted here that the Christian parents have had the child baptized, that the child is now in covenant relation with God, a member of the household of faith, of the family of the redeemed. That child is now to be treated as a son or a daughter of God. Upon the parents rests the responsibility of bringing up that
\quellseite{24}
child in the nurture and admonition of the Lord. In that ideal Christian home religion approaches the child first in the form of life and then in the form of instruction. The Spirit of Christ should permeate the atmosphere in which the child lives, the air which he breathes. In such a home the child will be instructed from the beginning. He will learn to pray as soon as he can speak. He will use his first lispings and stammerings in words of prayer. In the beautiful words of Luthardt:

“Let it not be objected that the child cannot understand the prayer. The way of education is by practice to understanding, and not understanding to practice. And the child will have a feeling and a presentiment of what it cannot understand. The world of heavenly things is not an incomprehensible region to the child, but the home of his spirit. The child will speak to his Father in heaven without needing much instruction as to who that Father is. It seems as though God were a well-known Friend of his heart. The child will love to pray. If mother forgets it, the child will not.”

In such a home there will be daily and diligent instruction in God’s truth, in order that the child may be sanctified through the truth. Bible pictures, Bible stories, Bible lessons, milk for babes, and gradually, little by little, the stronger meat of Bible doctrine.\randnotiz{What the child should be taught.} The Lutheran Graded Sunday School
\quellseite{25}
Series will prove a welcome guide and help to the parents. By and by the parts of Luther’s Small Catechism will be used in the home according to the intention and purpose of Luther “in the plain form in which they are to be taught by the head of the family.” As we have written elsewhere:

“The Apostles frequently speak of the Church that is in the house. By this they mean a home where the religion of our blessed Saviour permeates the whole atmosphere, where the Word of God dwells richly; where there are altars of prayer and closets of prayer—a home where Jesus is a daily, a well-known guest; where the children, baptized into Christ, are nourished with the milk of the Word, so that they grow thereby, increasing more and more, growing up unto Him who is the Head, even Christ. In such a home the Church is in the house and the household in the Church. Blessed home! Blessed children who have such parents! Blessed parents who have thus learned God’s way of grace!”\footnote{See Chapter VI., of “The Way of Salvation in the Lutheran Church.”}

\randnotiz{The Church and the child.} But such homes, alas! are only too rare. Had we more such homes, were they the rule instead of the exception, one of the Church’s greatest needs would be met. But since it is not so, since the great majority
\quellseite{26}
of children are reared in homes which are far from ideal, and in which they do not receive the spiritual nurture which they need, the Church must be the mother of her children, and, indeed, of all the neglected ones whom she can reach. She must furnish the needed teachers, and the needed opportunities for all who can be brought to learn what they need to be made wise unto salvation.

These teachers of the Church are called Catechists. The name is significant. Its derivation is interesting. Its Biblical use is instructive.\randnotiz{Derivation of Catechist.} It is a compound word.

The Greek noun is \emph{ἠχή}, which means sound, a spoken word, a word sounded back, an echo. The verb is \emph{ἠχέω}, to give back a sound, to echo, to answer. Classically it was used of the sounding down of the rushing water, of the sound of music falling from the ship on the sea. Then it came to signify the sounding down of a word or words of command or instruction from a superior to an inferior, from a teacher to a pupil. The preposition \emph{κατά} strengthens the meaning, bringing out more emphatically the back or return sound, the echo, the answer. Thus it came to mean instruction by word of mouth, familiar, conversational instruction, a free informal discussion between teacher and pupil.

In the New Testament Luke informs Theophilus (Luke 1:4), that he intends to give him a succinct and orderly account of those things which he had
\quellseite{27}
previously received by word of mouth (Περὶ ὧν κατηχήθης). See further the Greek in Luke 21:25, Acts 2:2; 18:25; 21:21; Rom. 2:18; 1 Cor. 14:19; Gal. 6:6, et al.

\randnotiz{What a Catechist is.} A Catechist, then, is a teacher who instructs his pupils, by word of mouth, in a free, familiar, conversational way in the elements of the Christian religion. His purpose is that he may make his pupils wise unto salvation. He desires that his pupils shall learn to know about their God and Saviour, their relation and duty to Him, in order that they may know Him, whom to know is life eternal.

Whether he have one pupil or many; whether he use the lecture system or the free and extemporaneous address, in either case giving opportunity for reciprocal questions and answers; or whether he use his own questions and let the pupils frame their own answers; whether his pupils be children or adults, or both, in every case he is a Catechist. As we shall see, all the above methods have been used and are still used. We do not here decide which is the best.

In all his work the good Catechist constantly endeavors to train and develop the mind, as well as to impart instruction. He desires to promote clear thinking as well as right knowing. He would teach both how to think and what to believe. The two
\quellseite{28}
should always go together. \randnotiz{The aim of the Catechist.} The good Catechist will put all his instruction and information in such simple, concrete, attractive and interesting form as will make it fit and acceptable to the pupils. His instruction should fit the mind, touch the heart, stir the conscience and move the will.

In our day the ordinary and regular Catechist is the pastor of the flock, or the missionary who would gather a flock. In his regular catechetical work he has to do with the children of the Church, and with children whom he would gather in and build up in the Church. As he thus faces and teaches his class he is doing a work second to no other in responsibility and blessedness. It is certainly a work heavy with responsibility. The quaint old Fuller used to say: “Every youth can preach, but he must be a man indeed who can profitably catechise.”

\randnotiz{The responsibility of the Catechist.} Let the Catechist ever bear in mind, as are the children of to-day so will be the men and women of to-morrow. Those who are children now are to fill the positions of power and trust in the family, the Church and the State of the future. How solemn the position. Souls are committed to his care; eternal interests depend on him. He is a watchman to warn, a shepherd to protect, to lead, to feed. The eternal and temporal weal or woe may depend on his catechising. Let him properly feed these lambs, and he will bind them to himself and to the Church
\quellseite{29}
as with hooks of steel. He will fence out heresy and soul-destroying error. He will save their bodies and their souls from the paths of the destroyer. No pastor should go before his class without having been first on his knees. All this will become more clear as we proceed.

Before we close this chapter, a few words on the blessedness of this work.

\randnotiz{Blessedness of this work.} What work could be more interesting than to instruct, influence, mold and develop these open, plastic and eager young hearts and minds! Truly “he who scatters the seed of divine truth in the heart of a child is training a plant for the paradise of God.” Let him be faithful to his trust and these children will be his crown of rejoicing in the day of the Lord. Here also it is true:

“’Tis not a work of small import
The pastor’s care demands;
But what might fill an angel’s heart
And filled a Saviour’s hands.

“To teach the way of Christ and Peace
It is a Christlike thing.”

A New England novelist, poet, journalist and jurist,\footnote{John Neal.} who had a clear insight into the illuminating and satisfactory teaching of the old Bible on child-nature and child-needs, writes thus seriously: “Even fathers and mothers look upon children
\quellseite{30}
with a strange misapprehension of their dignity. With the poet they are only the flowers and blossoms, the dewdrops or the playthings of earth. . . . Flowers! They are the flowers of the invisible world: indestructible, self-perpetuating flowers, with each a multitude of angels and evil spirits underneath its leaves, toiling and wrestling for dominion over it. Blossoms! They are the blossoms of another world, whose fruitage approaches angels and archangels. Or dewdrops! They are dewdrops that have their source, not in the chambers of earth, nor among the vapors of the sky, which the next breath of wind, or the next flash of sunshine may dry up forever, but among the everlasting fountains and inexhaustible reservoirs of mercy and love. Playthings! If the little creatures would but appear to us in their true shape for a moment, we should fall upon our faces before them, or grow pale with consternation, or fling them off with horror and loathing.”

Good words for the Catechist to ponder as he goes from his closet to his class.

\quellseite{31}
\chapter[Chapter II. The Catechist’s Activity]{Chapter II. The Catechist’s Activity. Its Relation to Church-Growth}

The Lutheran Church has many distinctive and distinguishing glories. Not the least of these glories is her system and custom of training her children and youth. In this she is peculiar. It is claimed, without hesitation, that no other Church has done and is doing so much for her children as our Evangelical Lutheran Church.

\randnotiz{The Lutheran Church and Christian nurture.} Wherever our Church is true to herself every pastor is a Catechist. For at least one term, if not for two terms of six or eight months, he meets the children who are of a proper age for his religious instruction and catechises them once or twice every week. He teaches them in systematic order the great foundation facts and truths of God’s word as they have been gathered and arranged in Luther’s Small Catechism. No other Church has laid so much stress on such instruction of her youth as our Church has. No other Church has a literature so rich and so helpful on this subject as has our Church. No other Church has been so generally painstaking and faithful as has ours. And we hesitate not to claim that no other Church has a body of young
\quellseite{32}
people so well informed and indoctrinated on the great facts and teachings of the Bible as has ours.

We are often reminded of an incident of some years ago. We were sitting in a sleeping car going from Fargo to Minneapolis. \randnotiz{An illustration.} A good woman from Fargo, a leader in the best social and church life of the city, an active member of the Episcopal Church, came across the car, sat down and said she wanted to ask a question. She wanted to know what we Lutherans do with our young people to make them so intelligent in religion. We asked her why she asked. She informed us that she had had a number of Lutheran girls work for her, mostly Scandinavians, and that with scarcely an exception she found them surprisingly at home in the contents and teaching of the Bible, so much so that they often had a better knowledge and a clearer understanding than herself who had grown gray in the society and service of her Church. We explained to her our system of Catechisation. She was grateful and expressed the wish that they might have it so in their Church.

\randnotiz{Another illustration.} Another instance: We were in Pittsburgh, soliciting for our Seminary. We called on a wealthy and benevolent Presbyterian lady, who had been a personal friend and helper of the Passavants. We told the story of our struggling school and what we hoped to accomplish in furnishing English-speaking ministers for the great army of children and youth
\quellseite{33}
of foreign, but Lutheran, parentage. She was interested, told us of the western institutions of her own Church, and how hard it seemed to get support for them. And then, as if soliloquizing, she said: “There is something wrong with the Presbyterian Church.” We asked her what she thought was wrong. She answered unhesitatingly: “Too much Christian Endeavor.” And then she went on to tell how, when she was young, she was catechised and taught to understand and love the teachings and practices of the Presbyterian Church, and how the youth of that Church knew why they were Presbyterians and loved their own Church, and how now in the Endeavor Societies they are impressed with the idea that one Church is as good as another, and that it makes no difference to which one they belong, etc. “I wish our young people were like those of your Church,” she said, “because they know why they are Lutherans.”

In addition to and in harmony with our scheme of catechisation we have our graded Sunday schools, our other Church schools, our Luther Leagues with their helpful topics and reading courses. This whole scheme of instruction is intended to make steadfast and intelligent believers out of our children and youth, who are then expected to be ready always to give an answer to every man that asketh a reason of the hope that is in them. And so we
\quellseite{34}
believe that we have a right to count our peculiar system of instruction one of the glories of our Church and to claim that the old Church of the Reformation does more for her children than any other Church does.

The question naturally arises, what are the results of this great and good system of Christian nurture? Do the results substantiate the claims made for it? Can it be known and estimated by its fruits? Are the fruits proportionate to the claims put forth?

These are fair questions and deserve a frank and fair answer.

We answer: Yes; we can show much excellent fruit from the system. \randnotiz{Good results.} Of the catechised and confirmed youth that we hold as loyal members of our congregations, we can claim that they will measure up with the members of any other Church. As to intelligence and devotion to their Church; as to firmness against alluring isms and schisms, head-heresies and heart-heresies, they will average above those of other churches. If they do not come up in activity and generosity it is because they are not properly encouraged and directed. On the whole, we have no reason to be ashamed of, but much ground for being proud of the loyal, confirmed youth of our churches. We can say to those who doubt us: “Here they are; they can speak for themselves.”
\quellseite{35}
\randnotiz{Failure in results.} But, on the other hand, we admit frankly and sadly that in quantity the fruits come short, by far, from what the system should lead us to expect. The church-records of our older congregations, in only too many instances, tell a sad, sad story of lapses and losses. And, stranger still, the saddest of all stories are told all too often by the records of those foreign-speaking congregations which claim to be the most strict and thorough in their catechisation. Some of the pastors of these congregations are frank enough to admit that if they can hold twenty-five per cent. of those whom they confirm they are well pleased. Others admit that they can count on holding only ten per cent. There are congregations that count their yearly confirmations by the scores and the hundreds, and yet are no stronger to-day than they were ten or twenty years ago. Some are even weaker. Some can explain this by the change in population and the removal of their people. But hundreds cannot rightfully give this explanation of their losses. It is not hard to find in a house-to-house canvass, in many districts of all our large cities, that in block after block a large proportion, if not a majority, of the residents have been confirmed in some neighboring Lutheran church and are now living without a thought of church attendance or connection. They are utterly estranged from their mother-church and seem to have neither
\quellseite{36}
scruple nor pang for their sad submersion in a godless world.

Large numbers of others have apostatized from the faith in which they were confirmed, are worshiping in and supporting churches of another faith and another spirit. The records of many of the aggressive non-Lutheran churches also contain some startling revelations. Many of their best members, their most liberal contributors, the very pillars of the congregations and of the denomination’s interests, were confirmed at Lutheran altars. The language problem may account for much of this, but for not nearly all.

When we, therefore, fearlessly face the facts and frankly admit them, we must say that the Lutheran Church has been catechising and confirming a large, an alarmingly large, proportion of her children into the world and into churches of a faith and spirit alien to that which they were taught in their course of catechisation.

\randnotiz{Can these failures be corrected?} Now we know that to calculate on holding all would be an unwarranted expectation. The deceitfulness and wickedness of the human heart will always prove too strong for many. The power of the world, the flesh and the devil will overcome large numbers. The attractions of false faiths, less humiliating and more boastful than their own, will always draw some away. There always have been,
\quellseite{37}
there always will be losses. But need there be as many as there are? Must our Church, filling in at the top, ever leak out so much of her best blood at the bottom? Is there no way of minimizing these distressing losses?

Let us suppose that we could reduce them to a necessary minimum. What if our Church should, in the coming generation, learn to reasonably hold her children? It requires no prophet’s ken to see that she would fast forge to the front. She has the material. She has the attractive doctrine. She has the open door. Opportunity such as no other church has, beckons her on. The vision is bright. The land and the time are big with promise. The one question is: Can she hold her children? Why has she failed to hold all she might have held? What has been wrong? Surely the system is not wrong. The wrong, then, must lie in the practice. There must be defects in the Catechists or in their method. If so, can we discover the wrong? Can it be righted? We believe the wrong can be located. We believe it can be corrected. It is with the earnest desire to help all our students and pastors to become good Catechists; good in consecrated purpose, good in theory, good in practice, with ever-increasing good fruit from their practice that we write this book.

\quellseite{38}
\chapter[Chapter III. Study of the Catechumen]{Chapter III. The Catechist’s Study of the Catechumen}

There can be no good pedagogics without a clear psychology. The most brilliant and most thoroughly educated teacher must fail if he does not understand his pupils.\randnotiz{Catechist should know child-psychology.} A large number of bright, well-equipped, good-intentioned and hardworking teachers fail because they do not know those whom they are trying to teach. When that great father of philosophy and pedagogy made “Know Thyself” his fundamental principle, he took for granted that to know one’s self is to know those whom one desires to teach. But Socrates was a teacher of adults. He had little or no concern for children. The greatest of all teachers fulfilled the prophecy that He should gather the lambs in His arms and carry them in His bosom. He loved little children. He took them in His arms and blessed them. He noted their playing in the market-place and drew lessons from it for the older ones. He set a child in the midst of His disciples, the future teachers of the Church, and made that child teach them some much-needed lessons. The first great charge He laid upon the heart of the restored Peter was “Feed my lambs.” His great commission for the whole Church for all
\quellseite{39}
ages is summed up in two words, viz., implanting life and nurturing life, baptizing and teaching. The greatest of the Apostles corrected a prevalent and false family psychology and pedagogy when he warned the fathers not to provoke their children to wrath, but to bring them up by and in proper nurture and admonition. Jesus understood child-psychology and wanted His disciples to understand it. Paul understood it, and wanted the fathers in Christian families to understand it. No one can be a good Catechist without it. Many Lutheran Catechists fail because they lack here.

The children in our catechetical classes range from twelve to sixteen years of age. This is the age of early adolescence.\randnotiz{Early adolescence.} The pastor must understand the psychology of this age. If he does not he is bound to blunder. His blunders may lead to the loss and ruin of some souls.

We note a few fundamentals that the good Catechist needs to know. We can give suggestions only. Every pastor should make a diligent study of child-psychology.

It is in the early adolescent age, which sometimes reaches to twenty, that there is the most rapid physical growth. Many shoot up like weeds. Their health is affected.\randnotiz{Period of physical weakness.} Specialists in Europe and America have concluded from wide and exhaustive experimentation that more than one-half of the children
\quellseite{40}
of this age are physically unwell. The mental development cannot keep up with the physical. We see at once how the teacher who is ignorant of this vital fact may wrong and alienate the pupil.

This is a restless age.\randnotiz{Period of restlessness.} It is next to impossible for the child of this age to sit still. And yet how many teachers and Catechists lose their temper, their time and the time of the class by repeatedly roaring, “Keep still; sit still; don’t move your hands and feet.” The boy now loves the gymnasium, the ball game, the snow-ball fight. The girl loves her roller or ice skates, her jumping rope, the toboggan or hillside slide, the tom-boy games and contests. It is a bad sign if they are quiet and demure at this age. Instead of fretting and fuming and despairing because of the turbulent restlessness of this age, the wise teacher will make allowance, have patience, be kind and sympathetic, and direct the energy instead of repressing it. A change in position is helpful. Call upon the restless one to stand up and answer. Let the whole class stand up and sing a stanza. Have an occasional five-minute recess. The most energetic of the boys and girls have, as a rule, the best stuff in them. The quiet boy who sits in the corner and pores over a book will never make a leader. Don’t pick him out for the ministry. He may become a dreamer of dreams, a book-worm, an author on obscure and abstract themes. But the world will never know that he has lived.

\quellseite{41}
It is at this age, also, that a continuous change is taking place in the sexual organism and functioning.\randnotiz{Period of sexual change.} This changing affects the physical health. More serious often is its influence on the psychical nature. It is a storm and stress period. The passions assert themselves. The whole emotional nature is disturbed. There is an upheaval of the whole being. The boy does not know what to make of himself. He is irritable, restless, suspicious. He imagines that father and mother do not appreciate him and do not treat him rightly. Sometimes he inclines to become morose, at other times he bursts the bonds and becomes wild. Traditions and authority do not count with him. He longs to get out of himself and out of his galling restraints. He dreams of running away and becoming a hero. It goes without saying that not all boys are affected in the same degree or in the same direction.

The girl has her troubles also. Her nature has much in common with that of the boy. But she also has experiences peculiar to her sex. She inclines to become dreamy; she also frets and worries. But she does not try to get away from herself. She becomes more and more self-centered; she loves to steal away by herself and brood and dream; she also feels misunderstood, unappreciated and neglected. She keeps a diary; she tells her diary her woes. What she cannot confide to the
\quellseite{42}
best of mothers, she discloses to the dear diary. She weeps over her fate and fancies herself the most forlorn girl in all the whole, wide world. Here, again, are different degrees and different directions. But the general peculiarities of the period are more or less manifest. This period may be overlapped by the former restless and active one, or that one may gradually merge into it. As these melancholic youth misunderstand their parents, so they are likely to form misconceptions of their church, their pastor and their Sunday school teacher. If their parents could only understand the psychology of this age they could greatly help them. As it is they often seriously wrong them. How easy for the pastor to wrong them if they are in his class at this critical age. And how much he can help them if he knows how.

Stretching through and more or less interwoven with the foregoing is a new emotional period. The dreaminess and brooding start longings for the remote and higher self. Up to this time the child has lived more or less outside of himself. He has taken his directions from without and followed them without question. Now he is more subjective. He lives in his own feelings. He mistrusts those who cannot or will not feel like himself. He turns away from parents and teachers. It is the chum age.\randnotiz{New emotional period.} A companion is sought out who feels like himself, who
\quellseite{43}
can enter into the misgivings and rebellions. Rivalries and jealousies spring up. Like-minded ones go together. They map out careers. They confide to each other their secrets. They air their contempt for those they do not like. Sometimes a boy and a girl become chums. They become child lovers. Others imitate or berate them. Many steal away and secretly read excitable tales, doubtful novels, impure books. It is a danger age. Parents and teachers need to watch and guard carefully the companions and the reading. Good company and pure diversions and elevating reading matter must be supplied. Parent and pastor must make special efforts to be companions and trusted friends to these uncertain and drifting ones. Let it not be forgotten that fifteen thousand children of the adolescent age were arrested in Chicago during the year just past.

If rightly understood, appreciated and utilized this is a hopeful and fruitful period for the Lutheran Catechist. As his pupils live so largely on the emotional plane, dry reasoning and curt commanding will fall flat. They will alienate rather than win. It is the emotional nature that is alert and hungry. It is open and receptive. Let the teacher appreciate this. Let him be warm-hearted, confidential and loving in his teaching.\randnotiz{Hopeful and fruitful period to Lutheran Catechist.} These pupils will respond beautifully when shown how the doctrine of sin corroborates the feelings of sin and guilt; how the teaching
\quellseite{44}
of repentance awakens an echo in hearts that have had indistinct feelings of dissatisfaction with self and sorrow for sin; how the portrayal of Christ’s love and sacrifice for sin and for sinners will call out a glad trust and a budding love; how the beauty of holiness will appeal to their slumbering ideals, and how the fact that in their baptism all this was given them in germ potency, will satisfy doubt and misgiving. Yes, our Lutheran system fits the child-psychology, while the evangelistic or revival system cruelly violates it.

It is at this impressionable age that the Catechist can greatly assist in the proper choice of a vocation. Here he can win candidates for the ministry of mercy and for the ministry of the Word. Let him properly show how the career of the former appeals to and satisfies the cravings of a woman’s heart and nature, while the latter is a heroic yet loving service of God and of humanity.

It is at this period, also, that the reasoning power begins to assert itself. It is the questioning and arguing period.\randnotiz{Questioning and arguing period.} A mere word of authority will no longer answer. The questioner—he is full of questions—wants to know why it is or should be so. Unhappy the teacher who gets out of patience with the questioner. A curt and cold reproof may raise doubts not easily allayed. Here patience must have her perfect work. Every sincere question

\quellseite{45}
should be commended. The frank telling of doubt should be encouraged. No pains should be spared to clear it up and satisfy the doubter. The time of the class dare not be sacrificed. Time must be taken after class.

It is in the catechetical age that memory is still at its best.\randnotiz{Memory still at its best.} This is the time to store the mind with precious truth which, even when not fully understood, will be working capital for the spiritual life of the future. The good Catechist should know the psychology of memory, and especially that of the child.

The imitative instinct is also strong at this age. It is the hero-worshiping age.\randnotiz{Hero-worshiping age.} The boy has his hero in the bigger, stronger or mentally superior boy. The girl in the superior girl. These are their models. These they praise and imitate. Others are disparaged, if not despised. Their weaknesses are mimicked and held up to ridicule. Happy is the pastor if he knows this and will win the admiration of these pupils. If he has peculiarities, mannerisms in conduct, dress or method of teaching, let him beware. He will be mimicked. But let him never show that he is annoyed or angered by this. If he does, the tendency will increase until he loses the respect of his class. Better laugh with the class at what he cannot help.

A more serious matter is the spirit and habit of irreverence
\quellseite{46}
peculiar to this age.\randnotiz{Irreverence peculiar to this age.} The would-be smart boy and the giggling girl are in evidence. They are a sore trial to the Catechist. They think everything is funny. They laugh at everything that pretends to be smart or witty. They are ready to laugh at everything and at nothing. There is nothing sacred in their eyes. The Catechist must know that this is a characteristic of that troublesome age. It is not meant to be as bad as it is. It is thoughtlessness more than intentional irreverence. It cannot be scolded away. Let the Catechist have patience. First win the confidence and love of the class. Then speak kindly but earnestly of the sacredness of God’s Word and of God’s house, and of these holy lessons of life and of love; show them that mirth and fun are all right in their place, but this is neither the place nor the time for them. By careful instruction the good Catechist can overcome it. In flagrant cases there must be some kind of discipline, as taking a back seat till after class and then a private interview.

On the whole, then, let the Catechist spare no pains to know his pupils. Let him know the psychology of the adolescent period. Let him verify and adapt his theories to the boys and girls before him. Let him learn to know them individually. No two are precisely alike. Each one must be studied; each one must be won. The peculiar interests and difficulties of each one should be discovered. The Catechist
\quellseite{47}
who would win will take a kindly interest in what interests each one. He will be attentive and helpful in the little and trifling troubles which seem so big at this age. Blessed the Catechist who can win the adolescent girl! More blessed he who can understand, win and hold the adolescent boy!
